\documentclass[11pt]{article}

\usepackage[margin=1in]{geometry}
\usepackage{amsmath, amssymb}
\usepackage{graphicx}
\usepackage{booktabs}
\usepackage{array}
\usepackage{float}
\usepackage{xcolor}
\usepackage{caption}
\usepackage{enumitem}
\usepackage{listings}
\usepackage[hidelinks]{hyperref}
\usepackage[backend=biber,style=numeric,natbib=true]{biblatex}
\addbibresource{references.bib}

\definecolor{codegray}{gray}{0.95}
\lstset{
  basicstyle=\ttfamily\footnotesize,
  backgroundcolor=\color{codegray},
  breaklines=true,
  frame=single,
  columns=fullflexible,
  keepspaces=true,
}

\newcommand{\fwd}{\ensuremath{e^{\mathrm{fwd}}}}
\newcommand{\rev}{\ensuremath{e^{\mathrm{rev}}}}
\newcommand{\comb}{\ensuremath{e^{\mathrm{comb}}}}
\newcommand{\score}{\operatorname{score}}

\newcommand{\gnote}[1]{{\color{blue} {Muhammad: #1}}}

\title{\textbf{Mechanistic Analysis of Adversarial Token Injection \\
in Neural Information Retrieval}}
\author{Muhammad Ghoummaid}
\date{\today}

\begin{document}
\maketitle

\begin{abstract}
Neural rerankers such as monoT5 can be manipulated by keyword-stuffing
attacks that inject relevance-like tokens into a document to inflate its
score without increasing its true relevance. Prior work shows such
attacks succeed, but not what happens inside the model when they do. We
address this gap in \texttt{castorini/monot5-base-msmarco} on TREC DL19
with a $7\times3\times5=105$-attack sweep and six causal and
representational analyses. Only 12 of 105 attacks raise the mean score,
and success depends on token identity and position rather than repetition
alone. Activation patching localizes the attack to \textbf{late-encoder
self-attention} (layers 10--11), the dominant causal component for all
105 attacks, which also contaminates the encoder's representation of the
\textbf{query} itself. This signal is built by a sparse set of
\textbf{18 of 144} encoder heads and read downstream by \textbf{33 of
288} decoder head-slots, 31 of them cross-attention (22 of which,
flagged before a later sample-size revision, were evaluated for
steering); at 20 of those 22 heads a per-head \textbf{linear direction}
captures a small but consistently steerable component of the attack
effect, bidirectionally. These results localize the vulnerability to a
specific, targetable part of the network, suggesting concrete
directions for lightweight defenses.
\end{abstract}

% =====================================================================
\section{Introduction}
% =====================================================================

Search systems rank documents against a query, and modern pipelines almost
always do it in two stages: a cheap retriever casts a wide net, then an
expensive cross-encoder reranker cleans up the top of the list:
\begin{center}
\verb|Retriever -> Top 1000 -> Cross-Encoder Reranker -> Top 10|
\end{center}
The reranker is where quality is won or lost, and it is also where a specific
kind of manipulation becomes possible. Neural rerankers will hand a high relevance
score to a document that isn't relevant at all, as long as it is sprinkled
with words that \emph{look} like relevance cues, a failure demonstrated for
sequence-to-sequence models like monoT5~\cite{parry2024adversarial}. It matters
because the fix is not obvious: a page about car parts should not outrank a
genuinely useful restaurant listing just because someone pasted the word
``relevant'' into it five times.

We already know these attacks work. What we don't know is \emph{why}: which
part of the model actually falls for the trick, and how the trick travels
from the injected words to the final score. This paper answers that, with
a controlled 105-attack sweep and six complementary tools: layer-level
activation patching, per-head patching and ablation, a linear mean-diff
steering intervention, query/document positional patching, a logit lens,
and a score-versus-rank validation.

The experiments reveal a consistent processing pattern, told here as a
sequence of experiments, each one opening the question the next one
answers: the attack effect is constructed primarily in late-encoder
self-attention, which also spreads into the query representation, and is
then read downstream mainly through decoder cross-attention.

This is a first step, not a full account. We leave a deeper investigation
of exactly where and how the attack lives in the model, and turning this
localization into a real defense, to future work, sketched in the
Conclusion.

% =====================================================================
\section{Related Work}
% =====================================================================

This work sits between adversarial neural IR and mechanistic
interpretability. We study monoT5, a sequence-to-sequence cross-encoder that
scores relevance through a generative objective~\citep{nogueira2020document},
because it is widely used and vulnerable to adversarial wording.

Prior attacks often assume access to a target query or model feedback,
whereas broader content-injection work shows that retrievers and rerankers
can also be manipulated by inserted queries, keywords, or unrelated
text~\citep{parry2024adversarial,tamber2025illusions}.
\citet{parry2024adversarial} establish a particularly weak threat model:
query-independent preemption, token-stuffing, and rewriting attacks requiring
no model-weight access, although they assume knowledge of the model's prompt.
They show that effectiveness depends on token identity, position, and
repetition, but their analysis remains at the model's \emph{output}: it
identifies which attacks work, not which internal computations produce the
score change.

Two recent studies apply causal interventions to ordinary relevance
computation. \citet{chen2024axiomatic} use activation patching to identify
attention heads associated with term-frequency behavior in the TAS-B dense
bi-encoder. \citet{lu2025pathway} trace a BM25-like circuit in a BERT-based
cross-encoder, where early and middle-layer matching heads feed later
relevance-scoring heads. We extend this work to an encoder-decoder reranker
under adversarial inputs: we trace how injected tokens alter encoder
representations, distinguish query- and document-position effects, identify
the decoder cross-attention heads that read the signal, and test whether the
effect can be manipulated through a low-dimensional steering intervention.

% =====================================================================
\section{Methodology}
\label{sec:method}
% =====================================================================

Before any of the experiments, we need three things fixed: how monoT5 turns a
(query, document) pair into a number, what an ``attack'' is, and how we read a
causal claim out of a patch. Everything after this section reuses this
machinery.

\subsection{monoT5 as a one-number relevance judge}

monoT5 treats ranking as a yes/no text problem: for each (query, document)
pair it builds the prompt
\texttt{"Query: \{query\} Document: \{passage\} Relevant:"}, runs the T5
encoder, takes a single decoder step, and looks at the logits at that first
decoder position. The relevance score is just the gap between the logits
for \texttt{true} and \texttt{false}:
\begin{equation}
\score \;=\; \operatorname{logit}(\texttt{true}) \;-\; \operatorname{logit}(\texttt{false}).
\label{eq:score}
\end{equation}
Both are single tokens, so this difference is exactly the model's log-odds
of \texttt{true} over \texttt{false}~\citep{nogueira2020document}: positive
means the model favors \texttt{true}, though that's a confidence gap, not
a calibrated probability of relevance. There is no generation: one
forward pass, one number.

\subsection{The attacks, and a control that isolates them}

An attack injects a trigger token into the document before scoring. We swept
every combination of:
\begin{itemize}[nosep]
  \item \textbf{7 tokens:} \texttt{relevant}, \texttt{information},
        \texttt{relevance}, \texttt{important}, \texttt{true}, \texttt{false},
        \texttt{bar};
  \item \textbf{3 positions:} \texttt{start} (prepended), \texttt{end}
        (appended), \texttt{random} (scattered inside the passage);
  \item \textbf{5 repetition counts:} $1$--$5$,
\end{itemize}
which gives $7\times3\times5=105$ token-injection configurations, which we
call attacks throughout as shorthand, though Section~\ref{sec:results-strength}
shows only 12 of 105 actually succeed. This extends the keyword-stuffing
side of \citet{parry2024adversarial}'s taxonomy across trigger identity,
position, and repetition count. We don't test their other two families:
multi-token \textbf{preemption}, which inserts a prompt-structure phrase
like \texttt{"Relevant: true"} rather than a single trigger word, and
\textbf{LLM rewriting}, which paraphrases or summarizes the whole passage.
Both change more of the passage than our padded control can cleanly align
against.

The trick to a clean causal claim is that comparison. For each pair we build
three inputs over the same passage:
\begin{description}[leftmargin=2em,style=nextline]
  \item[Original.] The untouched document (reference only; its length differs
  from the attacked input, so we don't patch against it).
  \item[Padded control.] The injection slots, wherever the trigger token
  sits, are filled with \texttt{pad\_token\_id} at
  \texttt{attention\_mask=0}, so the encoder can't see them, while every
  other passage token keeps the same absolute position as in the attacked
  input, the key move that holds sequence length and position fixed, so
  the only thing that differs from the attack is whether
  the injected words are actually there.
  \item[Attacked.] The real adversarial passage, all tokens visible
  (\texttt{attention\_mask=1}).
\end{description}

The quantity every experiment tries to explain is the \textbf{attack delta}:
\begin{equation}
\Delta \;=\; \score_{\text{attack}} - \score_{\text{control}}.
\label{eq:delta}
\end{equation}
Since control and attacked differ only in the injected words, $\Delta$ is the
causal effect of those words. Positive means the attack raised the score;
negative means it backfired. For the causal-patching analyses
specifically, we keep only examples with $\Delta>10^{-4}$, not just away
from zero, but positive, so those results characterize the pathway of
\emph{successful} score-inflation instances, including successes pulled
from configurations whose \emph{mean} effect is negative, not the average
behavior of every configuration.

\subsection{Reading causality out of a patch}
\label{sec:method-patching}

To find \emph{where} the effect lives, we use activation
patching~\citep{meng2022locating} across
monoT5's 12 encoder / 12 decoder layers ($d_{\text{model}}=768$) and five
component types per layer (encoder self-attention, encoder MLP, decoder
self-attention, decoder cross-attention, decoder MLP): swap one piece of
the network between a clean run and an attacked run, and see how much of
the attack's effect on the score follows the swap, asked in two
directions that each answer a different question. At this level we
replace a whole component's output at every sequence position; later
sections narrow that down to just query or document positions
(Section~\ref{sec:results-query}), and later still to one attention head
at a time (Sections~\ref{sec:results-encheads}
and~\ref{sec:results-heads}).

\textbf{Is this component enough? (forward / sufficiency.)} Run the clean
control, but at layer $L$, component $C$, paste in the attacked activation, and
see how far the score moves toward the attacked value, an operational
sufficiency test: can the attacked activation reproduce the effect with
the rest of the network left alone?
\begin{equation}
\fwd_{L,C} \;=\;
\frac{\score_{\text{patched}} - \score_{\text{control}}}
     {\score_{\text{attack}} - \score_{\text{control}}}.
\label{eq:fwd}
\end{equation}
$\fwd\approx0$: not enough on its own. $\fwd\approx1$: enough. $\fwd>1$:
amplified downstream.

\textbf{Is this component needed? (reverse / necessity.)} Run the attacked
input, but paste the clean activation into that one spot, and see how much of
the attack disappears, an operational necessity test: does removing just
this activation kill the effect?
\begin{equation}
\rev_{L,C} \;=\;
\frac{\score_{\text{attack}} - \score_{\text{patched}}}
     {\score_{\text{attack}} - \score_{\text{control}}}.
\label{eq:rev}
\end{equation}
$\rev\approx0$: others compensate, not needed. $\rev\approx1$: removing it here
kills the attack. $\rev<0$: an opposing mechanism.

A component only counts as being on the pathway if it clears both bars,
evidence it carries attack-related information rather than proof of a
single, formally necessary-and-sufficient path, so we take the smaller of
the two:
\begin{equation}
\text{combined}_{L,C} \;=\; \min(\fwd_{L,C}, \rev_{L,C}).
\label{eq:combined}
\end{equation}
Everything is normalized by the same $\Delta$ (Eq.~\ref{eq:delta}), so a value
of $1.0$ means ``as strong as the full unpatched attack effect,'' regardless
of the raw logit scale.

\subsection{Setup details}

Everything runs on monoT5-base (\texttt{castorini/monot5-base-msmarco}) in
a BM25 $\rightarrow$ monoT5 pipeline on TREC DL19, matching
\citet{parry2024adversarial}, random seed 42 throughout, with up to 500
(query, document) pairs per attack (487 shared across all 105). After the
initial sweep, we select \texttt{relevant\_start\_5} (token
\texttt{relevant}, prepended, repeated 5$\times$) as the
\textbf{canonical attack} used for every deep dive: it has the
second-highest mean effect (Table~\ref{tab:topbottom}) and uses the
literal relevance cue \texttt{relevant}; the main localization results
are additionally evaluated across all 105 configurations. The steering
direction (Section~\ref{sec:results-meandiff}) is estimated and evaluated
in-sample at two tiers, a breadth tier pooling up to 10 examples across
all 105 attacks and a depth tier of up to 100 examples on the canonical
attack, an initial representational diagnostic rather than evidence of
held-out generalization. Sample sizes for every other analysis are stated
where that analysis is introduced.

% =====================================================================
\section{Experiments and Results}
\label{sec:results}
% =====================================================================

\subsection{First: which attacks even work?}
\label{sec:results-strength}

Before chasing the mechanism, we first need to understand the connection
between an attack's token, position, and repetition count, and whether it
actually succeeds. Only \textbf{12 of the 105 attacks raised the mean score};
the other 93 \emph{lowered} it. These are configuration-level averages: a
configuration with a negative mean can still contain individual
query--document pairs where the injection raises the score, exactly the
successful-instance pool the patching analyses draw from
(Section~\ref{sec:method-patching}). The strongest configurations use the
tokens \texttt{relevant} and \texttt{information}, placed at passage
boundaries and repeated several times; the biggest losers were any token
scattered at random with many repetitions, and the averages hide how big
the swings get: individual pairs reach $+6.29$ under the strongest
attacks, so this isn't a gentle nudge for everyone, it's a large shove
for some.

\begin{table}[h]
\centering
\caption{Strongest and weakest attacks by mean attack delta $\bar\Delta$
(top/bottom 5 of 105).}
\label{tab:topbottom}
\small
\begin{tabular}{lr|lr}
\toprule
Strongest attack & $\bar\Delta$ & Weakest attack & $\bar\Delta$ \\
\midrule
\texttt{information\_start\_5} & $+0.381$ & \texttt{true\_random\_5}      & $-0.298$ \\
\texttt{relevant\_start\_5}    & $+0.364$ & \texttt{relevant\_random\_5}  & $-0.287$ \\
\texttt{relevant\_start\_4}    & $+0.286$ & \texttt{false\_random\_5}     & $-0.286$ \\
\texttt{information\_start\_4} & $+0.244$ & \texttt{relevance\_random\_5} & $-0.262$ \\
\texttt{information\_end\_5}   & $+0.184$ & \texttt{bar\_random\_5}       & $-0.226$ \\
\bottomrule
\end{tabular}
\end{table}

Three things drive whether an attack works. \textbf{Position.} Position
matters mainly through its interaction with token and repetition, not as
a simple effect on its own. Averaged over all tokens and repetitions,
\texttt{end} is almost neutral ($\approx-0.027$) and \texttt{start} is
mildly negative ($\approx-0.053$), while \texttt{random} is consistently
damaging ($\approx-0.156$): no random-position configuration has a
positive mean effect. Yet every one of the
12 attacks that worked sits at a passage boundary (6 \texttt{start}, 6
\texttt{end}): a boundary placement can support an attack when it is
paired with the right token and enough repetition, while scattering
tokens through the middle of the passage never does.

\textbf{Repetition.} It amplifies whatever sign the position already has:
\texttt{relevant\_start} goes from $-0.105$ to $+0.364$ as repetitions climb
from 1 to 5, while \texttt{relevant\_random} goes from $-0.101$ to $-0.287$
over the same range.

\textbf{Token.} Averaged over positions and repetitions, only
\texttt{information} ($\approx+0.018$) and \texttt{relevant}
($\approx-0.015$) land near zero; everything else is clearly negative
($-0.09$ to $-0.13$).

So repetition amplifies whichever sign position and token already set, but
it isn't sufficient on its own: the same extra tokens can raise the score
or lower it, depending on what's injected and where. That rules out a
simple count-only explanation, adding more tokens doesn't uniformly raise
the score, which is exactly what the mechanistic analyses that follow set
out to explain.

\subsection{A quick detour: how does score relate to rank?}
\label{sec:results-scorerank}

Every experiment below explains a change in \emph{score}, while users
experience a change in \emph{rank}. Across 52{,}478 pooled examples,
score and rank changes are strongly correlated
(Spearman $\rho=0.764$), with an even stronger correlation for already
competitive documents ($\rho=0.833$ vs.\ $0.754$ elsewhere). No
configuration raises the mean score while lowering the mean rank
(Appendix~\ref{sec:appendix-scorerank}). We therefore use score change
as a cheap, candidate-set-free, practical proxy for rank movement in
the remaining analyses.

\subsection{Where the attack effect is built: late-encoder self-attention}
\label{sec:results-encoder}

Now the real question: where does the attack-induced score change arise?
Decoder cross-attention is a natural first suspect because it directly
reads the encoder output before the decoder produces the relevance score.

Instead, the dominant component is \textbf{encoder self-attention}. It is
the top causal component for all 105 configurations. Seventy peak at
layer 11, 32 at layer 10, and only three at layer 0 or 9. Peak combined
importance ranges from $0.492$ to $0.972$, with a mean of $0.727$.
Figure~\ref{fig:aggcombined} shows the same thing averaged over the sweep:
encoder self-attention climbs from $\approx0.37$ at layer 9 to $0.65$--$0.69$ at
layers 10--11, towering over every decoder component and the encoder MLP.

\begin{figure}[h]
  \centering
  \includegraphics[width=0.7\textwidth]{outputs/attack_comparison/aggregate_combined_heatmap.png}
  \caption{Aggregate combined importance $=\min(\text{forward},\text{reverse})$,
  averaged over all 105 attacks. Rows are layers, columns components. The signal
  sits in \emph{encoder self-attention} at layers 9--11, not the decoder.}
  \label{fig:aggcombined}
\end{figure}

Averaged across layers, encoder self-attention also has the largest
combined effect ($0.290$), approximately twice decoder cross-attention
($0.146$); decoder self-attention and decoder MLP remain small ($0.035$
and $0.017$).

The effect does not switch on at one layer, but grows gradually with
depth, consistent with attack-related information accumulating through
the residual stream. This localization is stable across the sweep:
successful instances from both the 12 configurations with positive mean
effects and the 93 with negative mean effects peak in the same
late-encoder route (Appendix~\ref{sec:appendix-canonical-heatmap}).
Whatever the trigger word, the attack-related signal becomes concentrated
near the top of the encoder, where the injected tokens alter the joint
encoder representation. Decoder cross-attention carries a real but
secondary effect, consistent with it reading an encoder representation
that has already been modified.

That last point raises a question we hadn't planned to check. If the
encoder's self-attention is busy mixing the injected tokens into the
passage, is it careful enough to keep that contamination out of the
\emph{query}?

\subsection{The attack spreads into the query representation}
\label{sec:results-query}

Although every trigger is inserted in the document, encoder self-attention
can propagate its effect into the query positions. We test this by
patching only document positions, only query positions, or both, while
leaving the template positions fixed.

For the canonical attack, the query-only effect remains near zero through
the early layers and rises at layers 9--11, peaking at layer 11 in
encoder self-attention ($0.405$) and decoder cross-attention ($0.191$).
The negative early-layer values for \texttt{doc\_only} and \texttt{both}
are an artifact we traced to the frozen template positions, not a real
effect: they vanish once every position is patched together, reproducing
the whole-layer result exactly.

\begin{figure}[h]
  \centering
  \includegraphics[width=0.85\textwidth]{outputs/plots/effect_by_condition.png}
  \caption{Combined patching effect per layer, one line per condition
  (\texttt{doc\_only}, \texttt{query\_only}, \texttt{both}), encoder
  self-attention and decoder cross-attention side by side (canonical attack
  \texttt{relevant\_start\_5}, $n=100$). \texttt{query\_only} stays flat
  near zero until it rises cleanly at layers 9--11; \texttt{doc\_only} and
  \texttt{both} swing sharply negative at early-to-mid layers, the
  patching artifact discussed in the text.}
  \label{fig:bycondition}
\end{figure}

This pattern generalizes across the sweep: at encoder layer 11, the
query-only combined effect exceeds $0.2$ for all 105 configurations
(range $0.29$--$0.66$; Appendix~\ref{sec:appendix-queryacrossattacks}).
Thus, before decoder readout begins, the encoder has already written the
attack signal into both document and query positions. Independent
attention-mass and template-position checks are reported in
Appendices~\ref{sec:appendix-attnmass} and
\ref{sec:appendix-templatepositions}. Those checks also show that the
frozen template markers (\texttt{Query}, \texttt{Document}, and their
colons) carry some late-layer effect of their own, so the
\texttt{doc\_only} and \texttt{both} magnitudes above should not be read
as a complete decomposition of the attack effect.

``Encoder self-attention'' is really 12 heads
working in parallel at each layer, so that's still too coarse an answer to be
satisfying: is the effect spread evenly across all of them, or is it a
handful doing the work, with the same question waiting for us on the decoder
side once we get there?

\subsection{Who builds it: a sparse set of encoder self-attention heads}
\label{sec:results-encheads}

We next patch each of monoT5's $12\times12=144$ encoder self-attention
heads separately: each head writes an independent $64$-dimensional slice
of the layer's output, so it can be patched without touching the rest of
the residual stream. We use the same forward, reverse, and combined
measures as above.

It's just as sparse as we should have expected by now. Only \textbf{18 of
144} head-slots have a mean combined effect above $0.02$ across the attack
sweep, and the top ones are robust across the whole grid: \texttt{L11-H3}
and \texttt{L11-H4} clear it on \textbf{104 of 105} configurations,
\texttt{L10-H6} on 103, \texttt{L10-H2} on 102 (Figure~\ref{fig:encheadbox}).
The strongest encoder heads are divided between layers 10 and 11: the two
single largest individual effects, \texttt{L10-H0} ($0.182$) and
\texttt{L10-H2} ($0.161$), sit at layer 10 (full top-5 breakdown,
including forward/reverse, in Appendix~\ref{sec:appendix-enctable}).

\begin{figure}[h]
  \centering
  \includegraphics[width=0.7\textwidth]{outputs/plots/exp11_per_head_effect_across_attacks.png}
  \caption{Per-head combined effect distribution across all 105 attacks, top
  20 encoder heads by median. A tight box well above zero marks a head with
  a broadly consistent effect.}
  \label{fig:encheadbox}
\end{figure}

Single-head effects should not be summed to predict a multi-head
intervention: they come from different interventions and, in general,
different layers, so the effect of patching several heads together must
be evaluated jointly (Appendix~\ref{sec:appendix-encoderhead-behavior}).

Attention magnitude does not explain causal importance: only two layer-9
heads are both high-attention and causally important, while most of the
strongest heads in layers 10--11 have ordinary attention patterns.
Moreover, the trigger does not consistently act as a query-attention hub
across important heads. Thus, attention patterns alone do not reveal
which heads build the attack signal (full analysis in
Appendix~\ref{sec:appendix-encoderhead-behavior}).

The decoder's only job from here is to \emph{read} this now doubly-corrupted
representation, built by a sparse set of encoder heads and already leaking
into the query. Which of its heads actually do that?

\subsection{Who reads it: mostly decoder cross-attention heads}
\label{sec:results-heads}

We apply the same per-head patching analysis to the decoder's
$12\times12$ self-attention and $12\times12$ cross-attention heads, for
288 head-slots in total. We also run two ablations on the attack input:
\emph{zero} ablation sets the head's slice to zero, and \emph{mean}
ablation swaps in that same example's own matched padded-control
activation, a distribution-matched substitute, not an average across
control runs; the distinction turns out to matter a lot below.

The reading is remarkably sparse. Only \textbf{33 of 288 head-slots} (11.5\%)
clear a mean combined effect of $0.02$, and \textbf{31 of those 33 are
cross-attention} heads. Decoder self-attention is essentially inert (mean effect
$-0.00002$), while cross-attention's mean ($0.0040$) rises with depth and peaks
at layer 11 ($0.0249$), the same top-of-stack story, now on the decoder side.

\begin{figure}[h]
  \centering
  \includegraphics[width=0.85\textwidth]{outputs/plots/head_heatmap_grid_a.png}
  \caption{Mean combined effect per decoder head, over all 105 attacks ($n=100$
  each). Almost nothing in self-attention; a handful of late cross-attention
  cells carry the signal.}
  \label{fig:heatmap_a}
\end{figure}

The important heads show different response profiles.
\textbf{L11-X-H3} is near-universal: positive on 104 of 105 configurations
(median $0.069$, range $[-0.096,0.721]$), turning negative only on
\texttt{bar\_random\_4}. \textbf{L9-X-H2} reaches a large effect on some
configurations ($0.293$ on \texttt{true\_end\_4}) but goes inert
($<0.01$) on 25 of the 105: a defense that watches only one of these
would miss the other's failure mode entirely. Forward and reverse
patching agree closely across the decoder heads (mean
$|\fwd-\rev|=0.00034$), indicating that the same heads are selected by
both operational tests.

Here the choice of ablation is load-bearing, and it's a warning for
anyone building a defense on top of this. \textbf{Zero ablation
over-states and mis-attributes head importance}: zeroing
\texttt{L11-X-H3} raises both the attack score and the \emph{clean,
unattacked} score by nearly the same amount ($+0.437$ and $+0.450$), so
that head has a large general effect on scoring, not an attack-specific
one. \textbf{Mean ablation} gives the more honest picture: the same
head's attack-specific score drop is only $0.015$, small
but genuinely attack-specific, while the clean score barely moves
($0.0001$), and the same split holds for every top-5 head. The lesson
for any future head-ablation defense: measure with mean,
distribution-matched ablation, not zero, or you'll confuse ``this head
matters to the model'' with ``this head carries the attack'' (full
picture across all 288 heads, including the top-5 table, in
Appendix~\ref{sec:appendix-zeroablation}).

We now know where the signal is built, that it spreads to the query, and which
heads read it out. The obvious next question for anyone hoping to defend the
model: can we actually push on those heads?

\subsection{Can a linear direction steer the attack?}
\label{sec:results-meandiff}

Activation patching requires a matched donor run to pull a control
activation from, and mean ablation needs one too. We therefore ask a
sharper, more practical question: can a fixed direction, estimated in
advance, reproduce part of the same effect without access to a control
activation at inference time? Following the mean-difference
activation-steering approach~\citep{rimsky2024steering}, for each
flagged head (at layer $L$, component $C$, head $h$) we computed that
direction as the difference of class means over a pool of examples:
writing $a_i$ for the head's activation on example $i$,
\begin{equation}
\vec{d}_{L,C,h} \;=\;
\text{mean}_{i}\!\left[a^{\text{attack}}_{i}\right] -
\text{mean}_{i}\!\left[a^{\text{control}}_{i}\right],
\label{eq:meandiff}
\end{equation}
the average attacked activation minus the average control activation. We
then tested that single vector both ways, at the same intervention site
as per-head patching. \emph{Subtraction test:} on an attacked input, set
$a^{\text{mod}}=a^{\text{attack}}-s\,\vec{d}$ and see how far the score
walks back toward control. \emph{Addition test:} on a control input, set
$a^{\text{mod}}=a^{\text{control}}+s\,\vec{d}$ and see how far it walks
toward attack-like values. Scales $s\in\{0.5,1.0,1.5\}$. Of the 33 heads
identified in Section~\ref{sec:results-heads}, 22 were available for this
analysis: they were flagged by an earlier run of the head-patching sweep
at 10 examples per attack, before that sweep was later re-run at 100
examples per attack for statistical power (the 33-head figure reflects
the later run). Directions are fit and evaluated on the same pool, so
this is an in-sample feasibility test, not a generalization claim.

At \textbf{20 of the 22 flagged heads}, the subtraction test moves the
score back toward control, monotonically in scale: $1.5$ beats $1.0$
beats $0.5$ at every one of them, and none overshoots the baseline at
any scale. The best head recovers $6.8\%$ of the attack's score gap at
scale $1.5$ (Figure~\ref{fig:defensebyscale}), and it hasn't saturated
at the largest scale we tried. So a per-head linear direction captures a
small but consistently steerable component of the attack effect, not
the whole of it.

\begin{figure}[h]
  \centering
  \includegraphics[width=0.7\textwidth]{outputs/plots/defense_by_scale_grid_a.png}
  \caption{Mean score movement back toward the control baseline
  ($\Delta_{\text{toward control}}$) per flagged head, one line per scale,
  pooled over all 105 attacks. Heads sorted by scale-1.5 effectiveness.
  $1.5$ beats $1.0$ beats $0.5$ at every head except the last two
  (\texttt{L11-S-H3}, \texttt{L10-X-H0}), the only non-responders.}
  \label{fig:defensebyscale}
\end{figure}

The direction is also bidirectional: the addition test on a control
input pushes the score toward attack-like values, with nearly identical
head rankings under addition and subtraction (Pearson $r=0.999$, same
two non-responders both ways). This supports bidirectional control by
the direction, though the recovered fraction of the attack remains
small. Attack-specific directions are stronger than the pooled
direction, suggesting the shared geometry varies across configurations
(Appendix~\ref{sec:appendix-steering}). This is not yet a deployable
defense: the intervention is fit and tested in-sample, on inputs already
known to be attacked, not evaluated on held-out attacks.

Everything above is patching, an active intervention. As a final check, we
tested the causal story against an independent, non-causal
read-out~\citep{nostalgebraist2020logitlens} (a
logit lens): the injected token is never directly promoted near the top
of the vocabulary at any layer, on any of the 105 attacks, consistent
with an indirect effect rather than direct copying of the trigger (full
analysis in Appendix~\ref{sec:appendix-logitlens}).

% =====================================================================
\section{Conclusion}
% =====================================================================

We set out to explain \emph{why} monoT5 falls for keyword stuffing, not just
\emph{that} it does, and traced the effect through a 105-attack sweep and six
interpretability analyses. Attack success is token- and position-dependent
rather than a simple repetition effect: only 12 of 105 configurations raise
the mean score, all of them start/end-position and repeated, using one of
just two tokens (\texttt{relevant} or \texttt{information}), and score
change provides a strong proxy for the rank metric of
prior work ($\rho=0.76$). Across successful instances drawn from all 105
configurations, the attack effect is \emph{built} in late-encoder
self-attention (layers 10--11), the dominant causal component throughout
the sweep, concentrated in a sparse set of 18 of 144 encoder head-slots,
and at those same layers the encoder's own representation of the query
gets contaminated by the document. The decoder then \emph{reads} that
corrupted signal through an equally sparse set of 33 of 288 head-slots,
31 of them cross-attention, and at most of the heads evaluated, a
per-head linear direction steers the score both ways, though it recovers
at most $6.8\%$ of the attack's effect in-sample. A representational
control rules out direct copying or promotion of the injected token,
consistent with an indirect rather than direct mechanism. But no single
description covers the whole mechanism: attention mass explains only a
fraction of which encoder heads matter, and the attack token behaves like
an active hub only on average, not at every important head.
Overall, the score-inflation vulnerability lives in a specific,
targetable, but not monolithic part of the network, and future
robustness work can aim there rather than treating the model as a single
switch.

\paragraph{Future work.} Three directions matter most. First, trace the
missing link between the encoder and the decoder directly with path
patching, to find which encoder heads causally feed which decoder heads.
Second, an actual defense: combine several of the important heads and
test a joint intervention against held-out attacks and normal ranking
quality, not just the single-head, in-sample feasibility check we ran
here. Third, test generalization: repeat the same localization across
model sizes, architectures, datasets, and a broader set of attack
families, to see how much of this mechanism is specific to
monoT5-base.
 
\printbibliography

\appendix
\section{Appendix}

\subsection{Score-versus-rank agreement in full}
\label{sec:appendix-scorerank}

Section~\ref{sec:results-scorerank} states the headline correlation; this
is the full breakdown. Rank change is defined as
$\operatorname{rank}_{\mathrm{control}}-
\operatorname{rank}_{\mathrm{attack}}$, so positive values indicate upward
movement. We define ``competitive'' documents as those originally among
the query's top-100 BM25 candidates before any attack, as opposed to
documents injected from outside that candidate set.

The per-attack scatter of mean score change against
mean rank change is close to a single upward line
(Figure~\ref{fig:scorerank}), and the agreement holds where a skeptic would
expect it to break: with only a fixed 100-candidate list, ``improving the
rank'' of a document that starts buried way outside it is a fairly
artificial test, so one might expect the correlation to be an artifact of
that majority subgroup. Instead the correlation is \emph{stronger}, not
weaker, on the documents that were already genuinely competitive to begin
with ($\rho=0.833$ vs.\ $0.754$).

\begin{figure}[H]
  \centering
  \includegraphics[width=0.7\textwidth]{outputs/plots/score_vs_rank_scatter.png}
  \caption{One point per attack configuration (105 points): mean score
  change vs.\ mean rank change, coloured by success rate. Dashed lines
  split the plot into agreement/disagreement quadrants.}
  \label{fig:scorerank}
\end{figure}

Of the 105 attack configurations, 12 land in the both-improve corner and
89 in the both-worsen corner; only 4 disagree in sign, and the opposite
corner (score up but rank flat/down) is completely empty. No attack ever
raises the average score without also raising the average rank, and the
four disagreers all sit within $0.05$ of zero score change (score-neutral
attacks with a tiny net rank gain), not real contradictions. The rank
side even reproduces the strength ranking: \texttt{information\_start\_5}
and \texttt{relevant\_start\_5} are again the two strongest attacks.

\subsection{Per-layer, per-component effects for the canonical attack
(\texttt{relevant\_start\_5})}
\label{sec:appendix-canonical-heatmap}

Figure~\ref{fig:canonicalheatmap} shows the forward and reverse patching
effects underlying the discussion in
Section~\ref{sec:results-encoder} (\texttt{relevant\_start\_5}, $n=100$), at
the same per-layer, per-component granularity as the aggregate view in
Figure~\ref{fig:aggcombined}. Both directions climb gradually across layers
9--11 rather than switching on at a single layer, and agree closely at the
peak (encoder self-attention: $0.65$ forward, $0.81$ reverse, at layer 11),
confirming that this single attack reproduces the gradual, residual-stream
pattern described in the text rather than being an artifact of averaging
over the 105-attack sweep.

\begin{figure}[H]
  \centering
  \begin{minipage}{0.49\textwidth}
    \centering
    \includegraphics[width=\textwidth]{outputs/plots/forward_layer_component_heatmap.png}
  \end{minipage}\hfill
  \begin{minipage}{0.49\textwidth}
    \centering
    \includegraphics[width=\textwidth]{outputs/plots/reverse_layer_component_heatmap.png}
  \end{minipage}
  \caption{Forward (left) and reverse (right) patching effects per layer and
  component, canonical attack \texttt{relevant\_start\_5} ($n=100$). Rows are
  layers, columns components. Both directions rise gradually with depth and
  peak at encoder self-attention, layer 11.}
  \label{fig:canonicalheatmap}
\end{figure}

The same localization appears among successful instances drawn from both
the 12 configurations with positive mean effects and the 93 with negative
mean effects. Thus, configuration-level effectiveness changes the
prevalence and strength of successful instances, but not their dominant
localization.

\subsection{Query-only effect across the full attack sweep}
\label{sec:appendix-queryacrossattacks}

Section~\ref{sec:results-query} states the headline number; this is the
full picture. This isn't a quirk of the canonical attack
(\texttt{relevant\_start\_5}). The encoder-layer-11 query-only effect
exceeds $0.2$ on all 105 configurations (range $0.29$--$0.66$), at
exactly the layer range (9--11) where Section~\ref{sec:results-encoder}
put the overall effect (Figure~\ref{fig:queryacrossattacks}). By the time
the encoder finishes, the attack has altered not just the document's
representation but the query's, through the same late self-attention
layers.

\begin{figure}[H]
  \centering
  \includegraphics[width=0.85\textwidth]{outputs/plots/query_only_effect_across_attacks.png}
  \caption{Distribution of query-only combined patching effect across all 105
  attacks, per layer. Both regions show the same late-layer rise as the
  canonical attack (\texttt{relevant\_start\_5}).}
  \label{fig:queryacrossattacks}
\end{figure}

\subsection{Attention-mass corroboration for the query-contamination result}
\label{sec:appendix-attnmass}

Section~\ref{sec:results-query}'s query-only patching result is causal; this
is the independent, non-causal check that agrees with it. We measured how
much one span attends to another, normalized against a token-count-uniform
baseline (a value above 1 means more than uniform):
\begin{equation}
\text{normalized\_mass}(A\to B) \;=\;
\frac{\text{mean}_{i \in A}\big[\sum_{j \in B} \text{attn}[i,j]\big]}
     {|B| / \text{total\_seq\_len}}.
\label{eq:attnmass}
\end{equation}
Document-to-query mass shows an isolated spike at layer 9 ($1.42$, over
40\% above baseline), an anomaly against every neighbouring layer and
lined up exactly with where the query-only patching effect switches on.
The reverse direction shifts more subtly: query-to-document mass stays
below baseline everywhere but ticks up under attack at \emph{every} layer
(e.g.\ layer 8: $0.173\to0.194$), a small, universal increase in how much
the query looks at the injected content (Figure~\ref{fig:attnmass}).

\begin{figure}[H]
  \centering
  \includegraphics[width=0.85\textwidth]{outputs/plots/attention_mass_clean_vs_attack.png}
  \caption{Normalized attention mass per encoder layer, clean vs.\ attacked
  (canonical attack \texttt{relevant\_start\_5}, $n=100$). Left:
  query$\to$document. Right: document$\to$query. Dashed line: uniform
  baseline (1.0).}
  \label{fig:attnmass}
\end{figure}

\subsection{Are the frozen template positions really inert?}
\label{sec:appendix-templatepositions}

Section~\ref{sec:results-query}'s positional patching always keeps 9 short
template positions frozen at the control's own value (\texttt{"Query:"},
\texttt{"Document:"}, \texttt{"Relevant:"}, and the final \texttt{</s>}), so
that only the query and document spans themselves get patched. That
convenience assumes those positions carry none of the attack's signal
themselves. We tested the assumption directly: single-position causal
patching of each of the 7 named positions individually (the leading space
token folds into \texttt{"Query"} as a tokenization artifact, so there are
7 distinct positions to test, not 9), on the canonical attack ($n=100$)
and, unconditionally, across the full 105-attack grid ($n=10$/attack).

Figure~\ref{fig:templatepositions} shows the result. \texttt{"Query"},
\texttt{"Document"}, and the colon immediately after each are not inert:
patching just one of these single words moves the score by itself,
robustly, on the large majority of attacks. The other three,
\texttt{"Relevant"}, its colon, and \texttt{</s>}, behave exactly like the
sinks they were assumed to be: patching them barely moves the score at
all, even though \texttt{</s>} receives more attention than any other
position in the entire prompt.

\begin{figure}[H]
  \centering
  \includegraphics[width=0.7\textwidth]{outputs/template_tokens/plots/template_position_heatmap.png}
  \caption{Mean combined patching effect per encoder layer, each of the 7
  template positions patched individually (canonical attack, $n=100$).
  \texttt{Query}, \texttt{Document}, and their colons (top four rows) light
  up at layers 10--11; \texttt{Relevant}, its colon, and \texttt{</s>}
  (bottom three rows) stay flat throughout.}
  \label{fig:templatepositions}
\end{figure}

Summing the 7 positions' individual peak effects gives $0.211$, well under
the $0.484$ peak of the existing whole-span patch at the same layers: a
sub-additive relationship, where positions interact rather than
contributing independently. This is evidence that the excluded positions
carry real signal, not an additive correction to apply to
Section~\ref{sec:results-query}'s numbers. And the asymmetry,
\texttt{</s>} absorbing $40\times$ uniform attention while being the
least causally important position measured, is a recognizable signature
of a T5-style attention sink, expected behaviour for a model with no
dedicated beginning-of-sequence token to serve that role instead.

\subsection{Top encoder heads in full}
\label{sec:appendix-enctable}

Section~\ref{sec:results-encheads} states the top two; this is the full
top-5 breakdown on the canonical attack, including the forward/reverse
split.

\begin{table}[H]
\centering
\caption{Top 5 encoder heads by mean combined effect (canonical attack,
$n=100$).}
\label{tab:enctophead}
\small
\begin{tabular}{lrrr}
\toprule
Head & $\comb$ & $\fwd$ & $\rev$ \\
\midrule
L10-H0  & 0.182 & 0.200 & 0.222 \\
L10-H2  & 0.161 & 0.172 & 0.276 \\
L11-H3  & 0.121 & 0.127 & 0.165 \\
L11-H11 & 0.117 & 0.132 & 0.149 \\
L9-H6   & 0.087 & 0.218 & 0.242 \\
\bottomrule
\end{tabular}
\end{table}

\subsection{Encoder head behavior: non-additivity and attention analysis}
\label{sec:appendix-encoderhead-behavior}

Section~\ref{sec:results-encheads} states the two headline findings in
this subsection; this is the full picture.

\paragraph{Single-head effects should not be summed.} It is tempting to
compare the sum of the flagged heads' individual combined effects against
the whole-component patch from Section~\ref{sec:results-query}'s
experiment. These values should not be compared additively, though:
they come from different interventions (patching one head at a time
leaves the other 11 free to work normally, while patching the whole
component replaces every head's contribution simultaneously) and, for
individual heads, potentially different layers, so there is no reason to
expect them to add up. The practical lesson is the same regardless: if
you want to build a defense that touches several heads at once, you
can't predict what it will do from single-head numbers, a direct test
requires jointly patching the selected heads, which we leave to future
work.

\paragraph{Attention does not explain importance.} Does the attention
spike at layer 9 (Section~\ref{sec:results-query}) explain why these
heads matter? Mostly, no. Only 2 of the 12 heads at layer 9 are both
high-attention and causally important; two other layer-9 heads have even
more attention but no causal effect at all, and the pattern disappears
completely at layers 10 and 11, where most of the important heads
actually live: four of the five strongest heads there show ordinary,
unremarkable attention.

Does the attack token act like a magnet, pulling in query information
more than an ordinary word would? On average, yes: the important heads
show the attack token paying more attention to the query than a normal
document word does ($+0.142$ vs.\ $+0.110$). But that average hides a
split. At one head, \texttt{L9-H3}, the attack token reaches for the
query more than anywhere else we measured ($+2.67$), a textbook example
of the magnet idea. But at \texttt{L9-H6}, the strongest layer-9 head by
causal effect, the pattern is reversed: the attack token attends to the
query \emph{less} than an ordinary document token ($-1.57$).

So neither attention magnitude nor the magnet hypothesis explains
importance on its own: attention patterns alone are not a reliable guide
to which heads build the attack signal.

\subsection{Zero vs.\ mean ablation across all 288 decoder heads}
\label{sec:appendix-zeroablation}

Section~\ref{sec:results-heads} states the verdict for the top 5 heads;
this is the full picture, including the top-5 table.

\begin{table}[H]
\centering
\caption{Top 5 decoder heads by mean combined effect (105 attacks, $n=100$).
X = cross-attention. drop$_{\text{zero}}$ / drop$_{\text{mean}}$: mean score
drop from zero / mean ablation (score before minus score after; negative
means the score rose).}
\label{tab:tophead}
\small
\begin{tabular}{lrrrrr}
\toprule
Head & $\comb$ & $\fwd$ & $\rev$ & drop$_{\text{zero}}$ & drop$_{\text{mean}}$ \\
\midrule
L11-X-H3  & 0.0822 & 0.0833 & 0.0837 & $-$0.437 & 0.015 \\
L11-X-H6  & 0.0677 & 0.0682 & 0.0691 & $-$0.193 & 0.009 \\
L11-X-H11 & 0.0500 & 0.0505 & 0.0512 & $-$0.325 & 0.012 \\
L11-X-H0  & 0.0466 & 0.0469 & 0.0478 & $-$0.123 & 0.008 \\
L9-X-H2   & 0.0391 & 0.0414 & 0.0424 & $-$0.155 & 0.011 \\
\bottomrule
\end{tabular}
\end{table}

Swapping in the head's own matched control activation (mean ablation),
rather than a raw zero vector, drops the clean, unattacked score by
essentially nothing ($0.0001$) for \texttt{L11-X-H3}, versus zero
ablation's $0.450$-point rise, the gap that makes zero ablation
untrustworthy as an attack-specific diagnostic. Figure~\ref{fig:scatter}
shows this holds across all 288
head-slots, not just the top few: the important heads sit far off the
drop(zero)$=$drop(mean) diagonal, confirming that zero ablation's large
apparent effect is mostly a generic disruption to the model's ordinary
functioning, not evidence the head is attack-specific.

\begin{figure}[H]
  \centering
  \includegraphics[width=0.7\textwidth]{outputs/plots/zero_vs_mean_ablation_scatter.png}
  \caption{Mean score drop from zero ablation (x) vs.\ mean ablation (y), per
  head. Important heads sit far off the drop(zero)=drop(mean) diagonal.}
  \label{fig:scatter}
\end{figure}

\subsection{Steering direction: attack-specificity and direction norm}
\label{sec:appendix-steering}

Section~\ref{sec:results-meandiff} states the headline steering result;
this is the rest of the picture. Fitting the direction to a
\emph{single} attack, rather than pooling across all 105, is
$1.5$--$3.4\times$ more effective on that attack, while keeping nearly
the same head ranking ($r=0.879$): the shared, pooled direction is a
real but conservative summary of a geometry that varies somewhat by
attack, and we leave a full geometric explanation to future work.

Direction \emph{size} is a weak predictor of a head's causal importance:
Pearson $r=0.293$ ($n=288$) between a head's mean-diff direction norm and
its combined patching effect (Section~\ref{sec:results-heads}). The
single largest-norm head in the whole decoder, \texttt{L10-X-H0}
($\sim9\times$ the population mean), is one of only two heads where the
steering intervention never helps (Figure~\ref{fig:defensebyscale}). The
practical implication: pick heads to steer by causal combined effect,
not by how large their direction happens to be.

\subsection{The injected token is never directly promoted (logit lens)}
\label{sec:appendix-logitlens}

Section~\ref{sec:results-meandiff}'s closing paragraph summarizes this
check; this is the full analysis. The simplest alternative story to the
causal one is that the important decoder heads just shout the injected
word into the output. To test it, at each decoder layer we took the
cross-attention sub-layer's own contribution to the residual stream
(output minus input) and projected it straight through the unembedding,
skipping the final layer norm, and tracked where the injected token lands
in the vocabulary.

\begin{figure}[H]
  \centering
  \includegraphics[width=0.85\textwidth]{outputs/plots/attack_token_trajectory.png}
  \caption{Attack-token rank (log scale; 1 = top of vocabulary) and raw logit
  per decoder layer, clean/control/attack overlaid, canonical attack
  (\texttt{relevant\_start\_5}, $n=100$). The injected token never gets near
  the top at any layer.}
  \label{fig:trajectory}
\end{figure}

It never gets promoted. Across all 12 layers, all 12 flagged heads, and all
105 attacks, the injected token never comes near the top of the
$\sim$32{,}100-token vocabulary. The best rank ever seen is
$\approx$1{,}776, and that's on the \emph{clean} input at layer 1, not the
attack; at layer 11 the attack's rank ($8{,}741$) is only mildly better
than clean's ($\approx$12{,}000). The projected top-10 is $\geq$89\%
``other'' at every layer, with the attack bucket at exactly 0\%
throughout. That's a genuine null, and it doesn't contradict the
patching: it refines it. Cross-attention's contribution is causally
important for the final score, but that importance is realized through
many later layers of processing (including the very layer norm this
projection skips). A component can change what downstream layers compute
without its own isolated output pointing anywhere near the answer. The
effect is mediated, not direct.

\end{document}